\section{Remote Meetings (RQ1)}

Definition Meeting (Abgrenzung zu WS und Meetings -> Eigentlich das gleiche)

Was sind remote Meetings?

% ---------------------------------------------------------------------------------------------------------------------
\subsection{Meetings in Modern Work: Trends and Context}
\label{sec:meetings-in-modern-work}

The Doolde Report "State of Meetings 2023" finds that 40\% of meetings are scheduled 1–4 weeks in advance (55\%
in North America). In contrast, the Mirsooft report reveals that most meetings (57\%) are ad hoc, without
formal invitations. \cite{bobbyDoodleStateMeetings2026}

Welche Arten von Meetings gibt es?

% ---------------------------------------------------------------------------------------------------------------------
\subsection{Key challenges: Off-topic discussions, unequal participation, lack of structure}
\label{sec:key-challenges-in-meetings}
A clear issue with goal-oriented meetings is the absence or unclear communication of objectives. It seems
plausible that a significant number of meetings are scheduled without a defined goal or agenda (as discussed in section
\ref{sec:meetings-in-modern-work}).
A lack of structure in meetings is therefore one of the key challenges.
While an agenda can serve as the foundation for such structure, it alone does not guarantee a successful
meeting. Meeting goals are often dynamic and may need to be adjusted as the meeting progresses, due to the associated
uncertainties \cite[chap 4.1]{chenAreWeTrack2025}. Although the agenda can outline a path to achieving the
overall meeting goal \cite[chap 2.1]{scottMentalModelsMeeting2024}, without linked sub-goals, it guarantees
neither the success of the meeting nor a timely conclusion. \cite[chap 4.1]{chenAreWeTrack2025}
Discussions can nevertheless get out of hand if participants do no know what to priorities.
\cite[chap 4.1]{chenAreWeTrack2025}
An agenda was identified by many respondents (67\%) as a key element for successful meetings in Doodle’s State
of Meetings Report 2019. Only the setting of clear meeting objectives (72\%) was cited even more frequently.
\cite[p. 11]{DoodleStateMeetings2019}

Based on this, the following summary challenge statement was derived:

\begin{quote}
  \textbf{Challenge 1: Clear goal-setting and structure.}
  \label{chal:clear-goals-and-structure}
  Goal-oriented meetings frequently lack clearly communicated objectives and a reliable structure, so that
  discussions lose focus and the meeting neither succeeds nor concludes on time.
\end{quote}

Chen et al. noted, as a secondary finding, that individuals in higher hierarchical positions tended to
intervene more frequently and maintain focus on meeting objectives when a clear list of goals was provided.
Conversely, those in lower hierarchical positions felt less obligated to intervene, even when they observed
deviations from the main topic. \cite{chenAreWeTrack2025}
Doodle’s Stat of Meetings Report 2019 identified a correlation between salary and self-perceived
participation in meetings. Respondents with higher salaries were more likely to describe themselves as active
participants. \cite[p. 6]{DoodleStateMeetings2019}
This may stem from inexperience, hierarchical constraints, group dynamics, or a lack of error culture, where
dissent risks loss of face. Low participation can narrow solution spaces, yield monotonous outcomes, and
reduce meeting effectiveness.

Dysfunctional communication—such as complaining, criticizing, and digressing—reduces effectiveness.
\cite[130]{kauffeldMeetingsMatterEffects2012}

-> Unausgeglichen / geringe participation
-> Unvorhergesehen veränderungen im Meeting (flasch Annahmen, meetingziel ändert sich)
-> technische Probleme: Instablie internetverbindungen, Störungen etc. (nicht im scope dieser arbeit aber
dennoch relevant)

Graph der zeigt das sich diese challanges alle gegenseitig beiflussen könnnen. (z.B. wenn nicht nicht alle
  ansichten geührt wurde, kann es später dazu kommen das wenn man eigentlich eien entscheidung treffen möchte
(convergent thinking) wieder zurück in divergent htink fällt, was zu zirkuläeren diskussionen führen kann.)

Umgekert, kann ein zulanges verweilen in der coonvergnen zone dazuführen, dass kein entshciedung getroffen
wird und kein fortschirtt wahrgenommen wird.

temkultur und hierarchie beinfluss ebenfalls die offen tielnehmer ideen teilen und wer welche ideen wie bewertet.
